\documentclass[10pt,a4paper]{amsart}
\usepackage{amsmath,amssymb,amsthm,graphicx}
\usepackage{verbatim}
\newenvironment{reference}{\comment}{\endcomment}
%\newenvironment{reference}{}{}
\newenvironment{slogan}{\comment}{\endcomment}
\newenvironment{history}{\comment}{\endcomment}

% For commutative diagrams we use Xy-pic
\usepackage[all]{xy}

% We use 2cell for 2-commutative diagrams.
\xyoption{2cell}
% Theorem environments.
%
\theoremstyle{plain}
\newtheorem{theorem}[subsection]{Theorem}
\newtheorem{proposition}[subsection]{Proposition}
\newtheorem{lemma}[subsection]{Lemma}

\theoremstyle{definition}
\newtheorem{definition}[subsection]{Definition}
\newtheorem{example}[subsection]{Example}
\newtheorem{exercise}[subsection]{Exercise}
\newtheorem{situation}[subsection]{Situation}

\theoremstyle{remark}
\newtheorem{remark}[subsection]{Remark}
\newtheorem{remarks}[subsection]{Remarks}

\numberwithin{equation}{subsection}

% Macros
%
\def\lim{\mathop{\mathrm{lim}}\nolimits}
\def\colim{\mathop{\mathrm{colim}}\nolimits}
\def\Spec{\mathop{\mathrm{Spec}}}
\def\Hom{\mathop{\mathrm{Hom}}\nolimits}
\def\Ext{\mathop{\mathrm{Ext}}\nolimits}
\def\SheafHom{\mathop{\mathcal{H}\!\mathit{om}}\nolimits}
\def\SheafExt{\mathop{\mathcal{E}\!\mathit{xt}}\nolimits}
\def\Sch{\mathit{Sch}}
\def\Mor{\mathop{\mathrm{Mor}}\nolimits}
\def\Ob{\mathop{\mathrm{Ob}}\nolimits}
\def\Sh{\mathop{\mathit{Sh}}\nolimits}
\def\NL{\mathop{N\!L}\nolimits}
\def\CH{\mathop{\mathrm{CH}}\nolimits}
\def\proetale{{pro\text{-}\acute{e}tale}}
\def\etale{{\acute{e}tale}}
\def\QCoh{\mathit{QCoh}}
\def\Ker{\mathop{\mathrm{Ker}}}
\def\Im{\mathop{\mathrm{Im}}}
\def\Coker{\mathop{\mathrm{Coker}}}
\def\Coim{\mathop{\mathrm{Coim}}}

% Boxtimes
%
\DeclareMathSymbol{\boxtimes}{\mathbin}{AMSa}{"02}

%
% Macros for moduli stacks/spaces
%
\def\QCohstack{\mathcal{QC}\!\mathit{oh}}
\def\Cohstack{\mathcal{C}\!\mathit{oh}}
\def\Spacesstack{\mathcal{S}\!\mathit{paces}}
\def\Quotfunctor{\mathrm{Quot}}
\def\Hilbfunctor{\mathrm{Hilb}}
\def\Curvesstack{\mathcal{C}\!\mathit{urves}}
\def\Polarizedstack{\mathcal{P}\!\mathit{olarized}}
\def\Complexesstack{\mathcal{C}\!\mathit{omplexes}}
% \Pic is the operator that assigns to X its picard group, usage \Pic(X)
% \Picardstack_{X/B} denotes the Picard stack of X over B
% \Picardfunctor_{X/B} denotes the Picard functor of X over B
\def\Pic{\mathop{\mathrm{Pic}}\nolimits}
\def\Picardstack{\mathcal{P}\!\mathit{ic}}
\def\Picardfunctor{\mathrm{Pic}}
\def\Deformationcategory{\mathcal{D}\!\mathit{ef}}
\usepackage{cleveref}
\crefname{originalreference}{original reference}{original references}
\Crefname{originalreference}{Original reference}{Original references}
\makeatletter
\newlabel{situation-quasi-compact-objects}{{0937}{source}{}{sourcepointer-1}{}}
\newlabel{situation-quasi-compact-objects@cref}{{[originalreference][0][]0937}{[1][0][]source}{}{}{}}
\makeatother
\crefname{originalreference}{original reference}{original references}
\Crefname{originalreference}{Original reference}{Original references}
\makeatletter
\newlabel{item-enough}{{0938}{source}{}{sourcepointer-2}{}}
\newlabel{item-enough@cref}{{[originalreference][0][]0938}{[1][0][]source}{}{}{}}
\makeatother
\crefname{originalreference}{original reference}{original references}
\Crefname{originalreference}{Original reference}{Original references}
\makeatletter
\newlabel{item-enough-qc}{{0939}{source}{}{sourcepointer-3}{}}
\newlabel{item-enough-qc@cref}{{[originalreference][0][]0939}{[1][0][]source}{}{}{}}
\makeatother
\crefname{originalreference}{original reference}{original references}
\Crefname{originalreference}{Original reference}{Original references}
\makeatletter
\newlabel{item-enough-qc-qs}{{093A}{source}{}{sourcepointer-4}{}}
\newlabel{item-enough-qc-qs@cref}{{[originalreference][0][]093A}{[1][0][]source}{}{}{}}
\makeatother
\crefname{originalequation}{original equation}{original equations}
\Crefname{originalequation}{Original equation}{Original equations}
\makeatletter
\newlabel{equation-towards-constructible}{{093D}{source}{}{sourcepointer-5}{}}
\newlabel{equation-towards-constructible@cref}{{[originalequation][0][]093D}{[1][0][]source}{}{}{}}
\makeatother
\crefname{originalreference}{original reference}{original references}
\Crefname{originalreference}{Original reference}{Original references}
\makeatletter
\newlabel{homology-lemma-characterize-weak-serre-subcategory}{{0754}{source}{}{sourcepointer-6}{}}
\newlabel{homology-lemma-characterize-weak-serre-subcategory@cref}{{[originalreference][0][]0754}{[1][0][]source}{}{}{}}
\makeatother
\crefname{originallemma}{original lemma}{original lemmas}
\Crefname{originallemma}{Original lemma}{Original lemmas}
\makeatletter
\newlabel{lemma-cokernel-map-towards-constructibles}{{093E}{source}{}{sourcepointer-7}{}}
\newlabel{lemma-cokernel-map-towards-constructibles@cref}{{[originallemma][0][]093E}{[1][0][]source}{}{}{}}
\makeatother
\crefname{originallemma}{original lemma}{original lemmas}
\Crefname{originallemma}{Original lemma}{Original lemmas}
\makeatletter
\newlabel{lemma-extension-towards-constructibles}{{093G}{source}{}{sourcepointer-8}{}}
\newlabel{lemma-extension-towards-constructibles@cref}{{[originallemma][0][]093G}{[1][0][]source}{}{}{}}
\makeatother
\crefname{originallemma}{original lemma}{original lemmas}
\Crefname{originallemma}{Original lemma}{Original lemmas}
\makeatletter
\newlabel{lemma-covering-gives-surjection}{{0934}{source}{}{sourcepointer-9}{}}
\newlabel{lemma-covering-gives-surjection@cref}{{[originallemma][0][]0934}{[1][0][]source}{}{}{}}
\makeatother
\crefname{originallemma}{original lemma}{original lemmas}
\Crefname{originallemma}{Original lemma}{Original lemmas}
\makeatletter
\newlabel{lemma-change-presentation-towards-constructibles}{{093F}{source}{}{sourcepointer-10}{}}
\newlabel{lemma-change-presentation-towards-constructibles@cref}{{[originallemma][0][]093F}{[1][0][]source}{}{}{}}
\makeatother
\title{The modules admitting the standard cokernel presentation form a weak Serre subcategory under a hypothesis on kernels}
\author{The Stacks Project Authors}
\date{stacks-project:sites-modules.tex:local-snapshot-20261012; GNU FDL 1.2 or later}
\begin{document}
\maketitle
\noindent\textit{Editorial note.} All delivered mathematical text is the author's own, copied verbatim from the pinned Stacks Project chapter file sites-modules.tex; no statement or proof step was rewritten, reordered or removed.\par\medskip
\noindent\textit{Editorial note.} The counted span cites ten external items (one homology criterion, one situation, several lemmas and an equation), all recorded in assembly.external\_references with canonical tags as navigation pointers.\par\medskip
\noindent\textit{Allowed standard prerequisites (editorial).}\begin{itemize}
\item Criterion for a full subcategory to be a weak Serre subcategory. Stacks Project, Homology, Lemma homology-lemma-characterize-weak-serre-subcategory (tag 02N3). Reduces the verification to the kernel case.
\item Closure of the standard presentation class under cokernels and extensions. Stacks Project, Modules on Sites, Lemmas lemma-cokernel-map-towards-constructibles (tag 09BC) and lemma-extension-towards-constructibles (tag 09BD). Supplies the two closure results that the reduction uses.
\end{itemize}\medskip
\ifdefined\hypertarget\hypertarget{sourcepointer-1}{}\fi
\noindent\textit{Source reference.} Original reference 0937: sources/stacks-chow/tags.txt, source lines 11505--11505. Cross-chapter standard prerequisite cited by the author inside the counted span; it is used as a pointer and its own proof is not part of this item.\par\medskip
\ifdefined\hypertarget\hypertarget{sourcepointer-2}{}\fi
\noindent\textit{Source reference.} Original reference 0938: sources/stacks-chow/tags.txt, source lines 11506--11506. Cross-chapter standard prerequisite cited by the author inside the counted span; it is used as a pointer and its own proof is not part of this item.\par\medskip
\ifdefined\hypertarget\hypertarget{sourcepointer-3}{}\fi
\noindent\textit{Source reference.} Original reference 0939: sources/stacks-chow/tags.txt, source lines 11507--11507. Cross-chapter standard prerequisite cited by the author inside the counted span; it is used as a pointer and its own proof is not part of this item.\par\medskip
\ifdefined\hypertarget\hypertarget{sourcepointer-4}{}\fi
\noindent\textit{Source reference.} Original reference 093A: sources/stacks-chow/tags.txt, source lines 11508--11508. Cross-chapter standard prerequisite cited by the author inside the counted span; it is used as a pointer and its own proof is not part of this item.\par\medskip
\ifdefined\hypertarget\hypertarget{sourcepointer-5}{}\fi
\noindent\textit{Source reference.} Original equation 093D: sources/stacks-chow/tags.txt, source lines 11511--11511. Cross-chapter standard prerequisite cited by the author inside the counted span; it is used as a pointer and its own proof is not part of this item.\par\medskip
\ifdefined\hypertarget\hypertarget{sourcepointer-6}{}\fi
\noindent\textit{Source reference.} Original reference 0754: sources/stacks-chow/tags.txt, source lines 9118--9118. Cross-chapter standard prerequisite cited by the author inside the counted span; it is used as a pointer and its own proof is not part of this item.\par\medskip
\ifdefined\hypertarget\hypertarget{sourcepointer-7}{}\fi
\noindent\textit{Source reference.} Original lemma 093E: sources/stacks-chow/tags.txt, source lines 11512--11512. Cross-chapter standard prerequisite cited by the author inside the counted span; it is used as a pointer and its own proof is not part of this item.\par\medskip
\ifdefined\hypertarget\hypertarget{sourcepointer-8}{}\fi
\noindent\textit{Source reference.} Original lemma 093G: sources/stacks-chow/tags.txt, source lines 11514--11514. Cross-chapter standard prerequisite cited by the author inside the counted span; it is used as a pointer and its own proof is not part of this item.\par\medskip
\ifdefined\hypertarget\hypertarget{sourcepointer-9}{}\fi
\noindent\textit{Source reference.} Original lemma 0934: sources/stacks-chow/tags.txt, source lines 11502--11502. Cross-chapter standard prerequisite cited by the author inside the counted span; it is used as a pointer and its own proof is not part of this item.\par\medskip
\ifdefined\hypertarget\hypertarget{sourcepointer-10}{}\fi
\noindent\textit{Source reference.} Original lemma 093F: sources/stacks-chow/tags.txt, source lines 11513--11513. Cross-chapter standard prerequisite cited by the author inside the counted span; it is used as a pointer and its own proof is not part of this item.\par\medskip
\begin{lemma}
\label{lemma-towards-constructible-when-serre-subcategory}
In Situation \ref{situation-quasi-compact-objects} assume
(\ref{item-enough}), (\ref{item-enough-qc}), and (\ref{item-enough-qc-qs})
hold. Let $\mathcal{O}$ be a sheaf of rings.
Let $\mathcal{A} \subset \textit{Mod}(\mathcal{O})$ be the full
subcategory of modules isomorphic to a cokernel as in
(\ref{equation-towards-constructible}).
If the kernel of every map of $\mathcal{O}$-modules of the form
$$
\bigoplus\nolimits_{j = 1, \ldots, m} j_{V_j!}\mathcal{O}_{V_j}
\longrightarrow
\bigoplus\nolimits_{i = 1, \ldots, n} j_{U_i!}\mathcal{O}_{U_i}
$$
with $U_i$ and $V_j$ in $\mathcal{B}$, is in $\mathcal{A}$, then
$\mathcal{A}$ is weak Serre subcategory of $\textit{Mod}(\mathcal{O})$.
\end{lemma}

\begin{proof}
We will use the criterion of
Homology, Lemma \ref{homology-lemma-characterize-weak-serre-subcategory}.
By the results of
Lemmas \ref{lemma-cokernel-map-towards-constructibles} and
\ref{lemma-extension-towards-constructibles}
it suffices to see that the kernel of a map $\mathcal{F} \to \mathcal{G}$
between objects of $\mathcal{A}$ is in $\mathcal{A}$. To prove this
choose presentations
$$
\bigoplus A_{V_j} \to \bigoplus A_{U_i} \to \mathcal{F} \to 0
\quad\text{and}\quad
\bigoplus A_{T_j} \to \bigoplus A_{W_i} \to \mathcal{G} \to 0
$$
In this proof the direct sums are always finite, and
we write $A_U = j_{U!}\mathcal{O}_U$ for $U \in \mathcal{B}$.
Using Lemmas \ref{lemma-covering-gives-surjection} and
\ref{lemma-change-presentation-towards-constructibles}
and arguing as in the proof of
Lemma \ref{lemma-extension-towards-constructibles}
we may assume that the map $\mathcal{F} \to \mathcal{G}$
lifts to a map of presentations
$$
\xymatrix{
\bigoplus A_{V_j} \ar[r] \ar[d] &
\bigoplus A_{U_i} \ar[r] \ar[d] &
\mathcal{F} \ar[r] \ar[d] & 0 \\
\bigoplus A_{T_j} \ar[r] &
\bigoplus A_{W_i} \ar[r] &
\mathcal{G} \ar[r] & 0
}
$$
Then we see that
$$
\Ker(\mathcal{F} \to \mathcal{G}) =
\Coker\left(\bigoplus A_{V_j} \to
\Ker\left(
\bigoplus A_{T_j} \oplus \bigoplus A_{U_i} \to \bigoplus A_{W_i}\right)\right)
$$
and the lemma follows from the assumption and
Lemma \ref{lemma-cokernel-map-towards-constructibles}.
\end{proof}

\end{document}
